% Generated draft: requires mathematical and source review.
\documentclass[11pt]{article}
\usepackage[margin=25mm]{geometry}
\usepackage{fontspec}
\setmainfont{Latin Modern Roman}
\usepackage{amsmath,amssymb,mathtools,unicode-math}
\setmathfont{Latin Modern Math}
\usepackage{enumitem,needspace,xurl,hyperref}
\setlength{\parindent}{0pt}
\setlength{\parskip}{5pt}
\setlength{\emergencystretch}{3em}
\newcommand{\N}{\mathbb{N}}
\newcommand{\Z}{\mathbb{Z}}
\newcommand{\Q}{\mathbb{Q}}
\newcommand{\R}{\mathbb{R}}
\newcommand{\C}{\mathbb{C}}
\newcommand{\F}{\mathbb{F}}
\newcommand{\A}{\mathbb{A}}
\newcommand{\PP}{\mathbb{P}}
\newcommand{\E}{\mathbb{E}}
\newcommand{\eps}{\varepsilon}
\newcommand{\dd}{\,\mathrm{d}}
\begin{document}
\section*{UnsolvedMath ID: 30003830}
\textbf{Source problem code:} catalogue entry
\subsection{Definitions and mathematical statement}
The source asks whether every finite extension of a decidable field is decidable. It does not specify a logical language. Under the usual interpretation using the pure first-order language of rings \(\mathcal L_{\mathrm{ring}}=\{0,1,+,-,\cdot\}\), define
\[
\operatorname{Th}_{\mathrm{ring}}(F)=\{\varphi:\varphi\text{ is an }\mathcal L_{\mathrm{ring}}\text{-sentence and }F\models\varphi\}.
\]
A field \(F\) is called decidable in this interpretation if an algorithm, given any \(\mathcal L_{\mathrm{ring}}\)-sentence, determines whether it belongs to \(\operatorname{Th}_{\mathrm{ring}}(F)\). This concerns the complete theory, without requiring a computable presentation of \(F\). For an extension of fields \(K\subseteq L\), its degree is \([L:K]=\dim_K L\), and it is finite when this dimension is finite.

With this language qualification, the original question asks whether
\[
\forall\,K\subseteq L\text{ fields},\qquad
\bigl(\operatorname{Th}_{\mathrm{ring}}(K)\text{ is decidable}\ \land\ [L:K]<\infty\bigr)
\Longrightarrow
\operatorname{Th}_{\mathrm{ring}}(L)\text{ is decidable}.
\]
There is no restriction on characteristic in the supplied statement. The assertion concerns decidability for each individual extension; it does not specify a uniform algorithm for all extensions.

\textbf{Unresolved source qualification.} The registry labels the record open, but its supplied statement-verification note says that the official report answers this question negatively by constructing a decidable field with an undecidable finite extension. Under the interpretation above, that reported negative answer would have the form
\[
\exists\,K\subseteq L\text{ fields}:\quad
[L:K]<\infty,\qquad
\operatorname{Th}_{\mathrm{ring}}(K)\text{ is decidable},\qquad
\operatorname{Th}_{\mathrm{ring}}(L)\text{ is undecidable}.
\]
This is a transcription of the supplied verification claim, not an independently verified result. The record marks its statement status as \texttt{unrecoverable} and supplies no clean replacement statement. Consequently, the displayed universal implication is preserved as the original question, with conflicting status information explicitly retained; it is not presented as a verified open problem.
\subsection{Short English statement}
The original question asks whether a finite field extension always preserves decidability of the complete first-order theory, provisionally read in the pure language of rings. The supplied registry contains conflicting status information: its open label and dated literature assessment coexist with a verification note reporting a negative answer. This statement-only draft requires human mathematical and source review.
\subsection{Sources}
\begin{itemize}
\item Local UnsolvedMath registry record, permanent ID \(30003830\), title ``Decidability of Finite Field Extensions,'' exact source problem code \texttt{catalogue entry}; registry location \texttt{data/unsolvedmath/problems.json}. No individual proposer is supplied. The recorded proposed year is 2018.
\item Supplied report citation: \emph{Field Arithmetic (2019)}, Oberwolfach Reports, DOI 10.4171/owr/2018/25, pp.\ 1517--1577. \url{https://doi.org/10.4171/owr/2018/25}. The supplied background describes Oberwolfach Reports as documenting mathematical workshops and their open-problem discussions.
\item Supplied literature source labeled ``Oberwolfach Report 2019'': \url{https://ems.press/content/serial-article-files/46747}.
\item Supplied literature source labeled ``Anscombe--Fehm'': \url{https://doi.org/10.2140/ant.2016.10.665}. The record supplies no title or precise passage connecting this reference to the question.
\item Literature assessment checked at \texttt{2026-08-21}: this is supplied dated curation, not independently verified. It labels the status open and says that the cited list remains a frontier for complete first-order decidability, with only restricted fragments and related languages settled. The background explicitly calls this dated literature triage derived from the supplied OWR curation and requests independent verification of time-sensitive status and cited claims.
\item The supplied statement-verification note reports a negative answer in the official report and says that this extraction is not an unresolved problem. It states that no separate clean statement is promoted and that the original record remains publicly visible for provenance. The record's statement status is \texttt{unrecoverable}. These supplied qualifications conflict with its open status and literature assessment; no source was independently consulted for this draft.
\end{itemize}
\end{document}
